\documentclass[a6paper,10pt]{article}
%\usepackage[T1]{fontenc}
\usepackage[british]{babel}
\usepackage[utf8]{inputenc}
\usepackage{float, graphicx,amsmath,amssymb,amsfonts,cite,enumerate,tabularx,caption,subcaption}
\usepackage[final]{pdfpages}
\usepackage{wrapfig}
\usepackage[margin=0.3in]{geometry}
\usepackage{../sidspaltHack}
%\usepackage[framed,numbered,autolinebreaks,useliterate]{mcode}
\newcommand{\mel}[1]{\small\textbf{\textit{mel. #1 \\}}}


\setlength{\oddsidemargin}{-0.37in}
\setlength{\textwidth}{215pt}

\pagestyle{empty}

\begin{document}
\nysida{14}{1}
\noindent
\huge{Oo Visor till Fysiker}
\begin{center}
\Large o1. Konglig Fysiks Paradhymn \\ 
\mel{Katyuscha}
\end{center}
\small Här på festen stiger åter glammet.\\
Sången börjar, tentan bortglömd är.\\
Lyss min strupe Du plågas utav dammet,\\
frukta ej ty hjälpen är just här\\
ibland oss.\\
Höj pokalen, dess flöde känns som sammet.\\
Drick till det som Bacchi vapen lär.
\vspace{5pt}\\
Känn, Quristina, känn hur blodet hettar.\\
Du, o Osquar, orsak till det är.\\
Timmar skrider och dygdens bojor lättar.\\
Fest och glädje kärleksflamman när\\
men minns att\\
blott ej synen en hungrig kärlek mättar.\\
Drick till det som Venus' vapen lär.
\vspace{5pt}\\
FYSIKER, gasqueropen de har skallat,\\
likt musik från någon högre sfär.\\
Tentans piska för länge har oss vallat.\\
Trotsa den och studiernas misär,\\
med lärdom\\
från de makter som ytterst har oss kallat\\
Bacchus, Venus värdar hos oss är.\\
Och vänner,\\
Bacchi nektar ej Venus' flamma släcker. \\
SKÅL för det Fysiks skyddsgudar lär. 
\vspace{5pt}\\
\textit{Under sista versen står fysiker, men inga andra, upp.\\ "Quristina" och "o Osquar" kan bytas efter behag.}
\begin{flushright}
\textit{Dum-Dum 1977}
\end{flushright}

\nysida{14}{2}
\setlength{\oddsidemargin}{-0.47in}
\noindent
\begin{center}
\Large o2. Årskursernas hederssång \\ 
\end{center}
\vspace{-15pt}
\begin{table}[!h]
\begin{tabularx}{0.75\textwidth}{l l}
\textbf{\small Alla:} & \textbf{\small För det var i vår ungdoms fagraste vår, }\\
 &\textbf{\small vi drack varandra till och vi sade "gutår".}\\
 &\textbf{\small Och alla så dricka vi nu \textit{Foo} till,}\\
\small Foo: &\small Och \textit{Foo} säger inte nej därtill...
\end{tabularx}
\vspace{-20pt}
\end{table}
\begin{center}
\textit{Vid Foo insättes n-m,\\där n=aktuellt årtal, m=0, 1, 2...}
\end{center}
%\vspace{-10pt}
\begin{minipage}[t]{0.3\textwidth}
\small 35 Fux\\
42 Fibonacci\\
51 Filibuster\\
52 Förutseende\\
54 Fafner\\
55 Femfemman\\
56 Fjalar\\
57 Fimbul\\
58 Fatal\\
59 Fritiof\\
60 Fredag\\
61 Fulbia\\
62 Figaro\\
63 Fter\\
64 Fyran\\
65 Frukt\\
66 Födde\\
67 Fnark\\
68 Flirp\\
69 Fiasko\\
70 Flera\\
71 Funilla\\
72 Fillhouse \\
73 Fadäs\\
\end{minipage}
\begin{minipage}[t]{0.3\textwidth}
\small 74 Förutan\\
75 Finans\\
76 Filur\\
77 Fräck\\
78 Fyffes\\
79 Fnusk\\
80 Fuktig\\
81 Flykting\\
82 Fiqus\\
83 Fikon\\
84 Filmjölk\\
85 Fortsätt!\\
86 Flyter\\
87 Fnöske\\
88 Flora\\
89 Fagott\\
90 FemFel\\
91 Fjortis\\
92 Fido\\
93 Fromhet\\
94 Flaxa\\
95 Ferdinand\\
96 Felande\\ Länken\\
\end{minipage}
\begin{minipage}[t]{0.3\textwidth}
\small 97 Fängsel\\
98 Folkdans\\
99 Fulöl\\
00 Formfranska\\
01 Flykta!\\
02 Fusklapp\\
03 Fetvadd\\
04 FanFar\\
05 Flörtfrisk\\
06 Friskus\\
07 Fokus\\
08 Førspel\\
09 Falsklarm\\
10 Framfusig\\
11 Frappera\\
12 Fumla\\
13 Frågvis?\\
14 Folkvett\\
15 Fanclub\\
16 Fuskbygge\\
17 Förarlös\\
18 Flingsalt\\
19 Fasett\\
20 Fotnot
\end{minipage}


\nysida{14}{3}
\setlength{\oddsidemargin}{-0.37in}
\noindent
\begin{center}
\Large o3. De Brevitate Vitae (Gaudeamus) \\ 
\end{center}
$\|$: Gaudeamus igitur, iuvenes dum sumus! :$\|$\\
Post icundam iuventutem, post molestam senectutem\\
$\|$: nos habebit humus. :$\|$
\vspace{5pt}\\
$\|$: Ubi sunt, qui ante nos in mundo fuere?:$\|$\\
Vadite ad superos, transite ad inferos,\\
$\|$: ubi iam fuere!:$\|$
\vspace{5pt}\\
$\|$: Vita nostra brevis est, brevi finietur,:$\|$\\
venit mors velociter, rapit nos atrociter,\\
$\|$: nemini parcetur.:$\|$
\vspace{5pt}\\
$\|$: Vivat academia, vivant professores!:$\|$\\
Vivat membra quaelibet, vivant membran quaelibet\\
$\|$: semper sint in flore!:$\|$
\vspace{5pt}\\
$\|$: Vivant omnes virgines, graciles formosae!:$\|$\\
Vivant et mulieres, tenerae,\\
$\|$: amabiles bonae laboriosae!:$\|$
\vspace{5pt}\\
$\|$: Vivat et respublica et qui illam regit!:$\|$\\
Vivat nostra civitas, Maecenatum caritas,\\
$\|$: quae nos hic protegit!:$\|$
\vspace{5pt}\\
$\|$: Pereat tristitia, pereant osores!:$\|$\\
Pereat diabolus, quivis antiburschius,\\
$\|$: atque irrisores! :$\|$ 
\begin{flushright}
\textit{Biskop Stada 1267, C.W. Kindleben 1781}
\end{flushright}

\nysida{14}{4}
\setlength{\oddsidemargin}{-0.47in}
\noindent
\begin{center}
\Large o4. Överföhssång \\ 
\mel{Kungssången}
\end{center}
\vspace{-15pt}
\begin{figure}[!h]
\begin{minipage}{0.6\linewidth}
\small Ur teknologens djup en gång\\
en samfälld och en enkel sång\\
som går till Överföhs fram
\vspace{5pt}\\
Var henne trofast och hennes ätt\\
Gör piskan i hennes näve lätt\\
Och all din tro till henne sätt\\
Du schlemm av ingen rang
\end{minipage}
\begin{minipage}{0.38\linewidth}
\includegraphics[width=\textwidth]{loggan.png}
\end{minipage}
\end{figure}
\vspace{-10pt}
\begin{flushright}
\textit{Skriven till Ettans Fest 1990}
\end{flushright}
\vspace{10pt}
\begin{center}
\Large o5. Studentsången \\ 
\footnotesize\textit{Noter till blandad kör finns i notkapitlet.}
\end{center}
Sjung om studentens lyckliga dag!\\
Låtom oss fröjdas i ungdomens vår!\\
Än klappar hjärtat med friska slag,\\
och den ljusnande framtid är vår.
\vspace{5pt}\\
Inga stormar än i våra sinnen bo,\\
hoppet är vår vän och vi dess löften tro\\
när vi knyta förbund i en lund.\\
Där de härliga lagrarna gro,\\
Där de härliga lagrarna gro.\\
Hurra!
\vspace{5pt}\\
Svea vår moder hugstor och skön\\
manar till bragd som i fornstora dar,\\
vinkar med segerns och ärans lön,\\
men den skörd utan strid man ej tar.

\nysida{14}{6}
\setlength{\oddsidemargin}{-0.37in}
\noindent
Aldrig slockne då känslans rena brand,\\
aldrig brista må trohets helga brand,\\
så i gyllene frid som i strid.\\
Liv och blod för vårt fädernesland!\\
Liv och blod för vårt fädernesland!
\begin{flushright}
\textit{Herman Sätherberg}
\end{flushright}
\vspace{20pt} 


\begin{center}
\Large o6. När Fumla blev fem \\
\mel{När vi två blir en}
\end{center}
\textit{Ettan:} \\
Det verkade lätt, att gå på KTH. \\
Jag var den med alla rätt, som ingen rådde på, \\
men det é så svårt å hinna med allt, \\
jag menar, plugga non-stop é inte så ballt. \\
Jag trodde jag var smart, \\
men teknisk fysik, va inte självklart.
\vspace{5pt} \\
\textit{Alla:} \\
Så vi dränker vår sorg med en fest. \\
Hör nu glasen slå! \\
Ty vårt valmanifest \\
är att shotta lite stroh. \\
Så lyft era glas, nu är det kalas \\
för termon ja ja, den gick väl bra. \\
Så vi dränker vår sorg med en fest, \\
nu börjar år 2.
\newpage
\setlength{\oddsidemargin}{-0.47in}
\noindent
\textit{Tvåan:} \\
Det verkade fett, att gå i årskurs två. \\
Om vi grejade år ett, så borde detta gå, \\
men att gasqua i kons är lite för kul, \\
så paniken slog till, lagom till jul. \\
Vi tar bäsken med en min, \\
och tentan får jag, klara nästa termin.
\vspace{5pt} \\
\textit{Alla:} \\
Så vi dränker vår sorg med en fest. \\
Hör nu glasen slå! \\
Ty vårt valmanifest \\
är att shotta lite stroh. \\
Så lyft era glas, nu är det kalas \\
för diffen ja ja, den gick väl bra. \\
Så vi dränker vår sorg med en fest, \\
nu börjar år 3.
\vspace{5pt} \\
\textit{Trean:} \\
Det verkade hett, att ta ett år med sol, \\
få leva mer jet set, å vara lite cool. \\
Men det é svårt å plocka poäng \\
när man hellre festar loss i en simbassäng \\
sen vare plötsligt vår, \\
med fasta och kex, jobbet består. 
\vspace{5pt} \\
\textit{Alla:} \\
Så vi dränker vår sorg med en fest. \\
Hör nu glasen slå! \\
Ty vårt valmanifest \\
är att shotta lite stroh. \\
Så lyft era glas, nu är det kalas \\
för kvanten ja ja, den gick väl bra. \\
Så vi dränker vår sorg med en fest, \\
nu kommer vi hem.
\newpage
\setlength{\oddsidemargin}{-0.37in}
\noindent
\textit{Fyran:} \\
Det verkade rätt, att grunda med komplex. \\
Ja nu blir mastern lätt, men vårat hopp det släcks, \\
för det é så svårt å göra inlupp \\
när man aldrig lärt sig jobba i grupp. \\
Det blir mest en massa prat, \\
å då blir det svårt, inga toppresultat. 
\vspace{5pt} \\
\textit{Alla:} \\
Så vi dränker vår sorg med en fest. \\
Hör nu glasen slå! \\
Ty vårt valmanifest \\
är att shotta lite stroh. \\
Så lyft era glas, nu är det kalas \\
för mastern ja ja, den går väl bra. \\
Så vi dränker vår sorg med en fest, \\
nu börjar år 5.
\vspace{5pt} \\
\textit{Femman:} \\
Hur har detta skett, alla åren rullar på. \\
Slipat CV, skötsamt sätt, så ska jag jobbet nå. \\
För nu är det dags, att tjäna sin deg \\
eller bli doktorand för å slippa kneg. \\
Måste lämna inom kort,  \\
så vi filar på vår exjobbsrapport.
\vspace{5pt} \\
\textit{Alla:} \\
Så vi dränker vår sorg med en fest. \\
Hör nu glasen slå! \\
Ty vårt valmanifest \\
är att shotta lite stroh. \\
Så lyft era glas, nu är det kalas \\
för examen ja ja, den blir nog bra. \\
Så vi dränker vår sorg med en fest, \\
å nu går vi ut.
\newpage
\setlength{\oddsidemargin}{-0.47in}
\noindent
\textit{Seniorer:} \\
Det var ju så lätt, att gå på KTH, \\
Har egen bostadsrätt, har nått en ny nivå. \\
De som pluggar nu, de har ingen koll, \\
jag menar dagens ungdom, de fattar noll. \\
Men en sak lever kvar,  \\
ingen har glömt, hur bra vi var.
\vspace{5pt} \\
\textit{Alla:} \\
När vi dränkte vår sorg med en fest, \\
å hörde glasen slå! \\
Ty vårt valmanifest \\
var att shotta MASSA stroh! \\
Så lyft era glas, nu är det kalas \\
för livet ja ja, vi har det bra. \\
Fast vi saknar att vara på fest \\
för vi har gått ut.

\begin{flushright}
\textit{F-12 Fumla \\ Jubilariegyckel, Ettans fest 2017}
\end{flushright}


\nysida{14}{$\infty$}
\setlength{\oddsidemargin}{-0.37in}
\noindent
\vspace{20pt} 
\begin{center}
\Large o$\infty$. O gamla klang och jubeltid \\ 
\small \textit{Noter finns i notkapitlet}
\end{center}
O gamla klang och jubeltid,\\
ditt minne skall förbliva.\\
Och än åt livets bistra strid\\
ett rosigt skimmer giva.\\
Snart tystnar allt vårt yra skämt,\\
vår sång blir stum, vårt skratt förstämt.\\
O, jerum, jerum, jerum!\\
O, quae mutatio rerum!
\vspace{5pt}\\
Var äro de som kunde allt,\\
blott ej sin ära svika.\\
Som voro folk av äkta halt,\\
och världens furstar lika.\\
De drogo bort från vin och sång,\\
till vardagslivets tråk och tvång.\\
O, jerum...
\vspace{5pt}\\
\textit{Filosofer:} Den ene vetenskap och vett\\
in i scholares mänger.\\
\textit{Jurister:} Den andre i sitt anlets svett\\
på paragrafer vränger.\\
\textit{Teologer:} En plåstrar själen som är skral.\\
\textit{Medicinare:} En lappar hop dess trasiga fodral.\\
O, jerum...

\newpage
\setlength{\oddsidemargin}{-0.47in}
\noindent
Men hjärtat i en sann student\\
kan ingen tid förfrysa.\\
Den glädjeeld som där hen tänt,\\
hens hela liv skall lysa.\\
Det gamla skalet brustit har,\\
men \textbf{kär}nan finnes frisk dock kvar.\\
Och vad hen än skall mista,\\
den skall dock aldrig brista.
\vspace{5pt}\\
Så sluten bröder fast vår krets\\
till glädjens värn och ära.\\
Trots allt vi tryggt och väl tillfreds,\\
vår vänskap trohet svära.\\
Lyft bägarn högt och klinga vän,\\
de gamla gudar leva än.\\
Bland skålar och pokaler!\\
Bland skålar och pokaler! 
\begin{center}
\textit{På fina fester dunkar en handen i bordet endast vid första stavelsen av "kärnan" i fjärde versen, inte annars.}
\end{center}
\begin{flushright}
\textit{Ur "Alt Heidelberg - R$\ddot{u}$ckblick eines alten Burschen"\\
Svensk text: August Lindh}
\end{flushright}

\end{document}
